\documentclass[11pt,a4paper]{article}

\usepackage[margin=1in]{geometry}
\usepackage[T1]{fontenc}
\usepackage[utf8]{inputenc}
\usepackage{lmodern}
\usepackage{booktabs}
\usepackage{longtable}
\usepackage{array}
\usepackage{tabularx}
\usepackage{hyperref}
\usepackage{enumitem}
\usepackage{xcolor}

\hypersetup{
    colorlinks=true,
    linkcolor=blue,
    urlcolor=blue,
    pdftitle={Model Progress, Bounds, and Run Results},
    pdfauthor={Unified VSD Model}
}

\newcolumntype{Y}{>{\raggedright\arraybackslash}X}

\title{Model Progress, Bounds, and Run Results}
\author{Unified VSD Pre-Surgery Calibration Workflow}
\date{2026-05-12}

\begin{document}

\maketitle

\section{Purpose}

This note records what has been done so far to make the model more
scientifically defensible, especially:

\begin{itemize}[leftmargin=1.5em]
    \item where the clinical data came from
    \item how the calibration bounds (\texttt{lb} and \texttt{ub}) are defined
    \item what changes were made to improve the model
    \item what the most important runs have shown so far
\end{itemize}

The goal is to keep one readable progress record that explains both the
scientific reasoning and the operational results.

\section{Data Sources Used So Far}

\subsection{Reyna}

Primary sources:

\begin{itemize}[leftmargin=1.5em]
    \item \texttt{Filled Protokol\_VSD Pre\_Reyna (2).pdf}
    \item \texttt{docs/clinical\_data\_dictionary.md}
    \item \texttt{docs/pak\_dipo\_co\_objective\_memo.md}
    \item \texttt{docs/reyna\_pak\_dipo\_systemic\_flow\_revision.md}
\end{itemize}

Key decisions:

\begin{itemize}[leftmargin=1.5em]
    \item Active anthropometry uses the raw protocol values:
    \begin{verbatim}
weight_kg = 13.4
height_cm = 95.0
BSA = 0.588 m^2
    \end{verbatim}
    \item Clinical \texttt{CO\_Lmin = 3.423} is compared to model
    \texttt{Qs\_Lmin}, not \texttt{LVCO\_Lmin}
    \item Pulmonary calibration is intentionally reduced to \texttt{PAP\_mean} only
    \item \texttt{PAP\_sys\_mmHg} and \texttt{PAP\_dia\_mmHg} are set to
    \texttt{NaN} in the active Reyna file
\end{itemize}

Why:

\begin{itemize}[leftmargin=1.5em]
    \item This keeps Reyna strictly patient-specific
    \item It avoids spending optimization leverage on pulmonary waveform shape
    that is less trustworthy than the mean pressure
    \item It follows Pak Dipo's advice that the main remaining issue is
    systemic preload-flow-volume consistency
\end{itemize}

\subsection{Razka}

Primary sources:

\begin{itemize}[leftmargin=1.5em]
    \item \texttt{config/patient\_profile\_Razka.m}
    \item procedure-log style catheter values embedded in the profile comments
\end{itemize}

Key characteristics:

\begin{itemize}[leftmargin=1.5em]
    \item Real patient data
    \item Good catheter pressure and \texttt{QpQs} information
    \item No measured \texttt{CO\_Lmin}
    \item No ventricular volume block
\end{itemize}

Why it matters:

\begin{itemize}[leftmargin=1.5em]
    \item Razka is useful for testing the sparse real-data path
    \item Razka is not the right case to test the Reyna CO-volume cleanup directly
\end{itemize}

\subsection{Profile A and Profile B}

Primary sources:

\begin{itemize}[leftmargin=1.5em]
    \item \texttt{config/patient\_profile\_A.m}
    \item \texttt{config/patient\_profile\_B.m}
    \item synthetic benchmark comments embedded in those files
\end{itemize}

Key characteristics:

\begin{itemize}[leftmargin=1.5em]
    \item These are synthetic benchmark profiles, not real patient cases
    \item They are useful for stress-testing the pipeline
    \item They should not be treated as scientific validation equal to Reyna or Razka
\end{itemize}

Important caution:

\begin{itemize}[leftmargin=1.5em]
    \item Profile A in particular contains internal tension between stated
    \texttt{MAP}, \texttt{Qs}, \texttt{SVR}, and volume targets
    \item Therefore a poor fit on A may reflect benchmark inconsistency rather
    than a true modeling failure
\end{itemize}

\section{What We Changed to Make the Model Better}

\subsection{Locked the CO definition to the correct comparator}

We changed the interpretation of clinical cardiac output so that:

\begin{verbatim}
clinical CO target -> model Qs_Lmin
not -> model LVCO_Lmin
\end{verbatim}

Reason:

\begin{itemize}[leftmargin=1.5em]
    \item In unrepaired VSD, LV output includes recirculated pulmonary/shunt flow
    \item Effective systemic delivery is represented by \texttt{Qs}
    \item Comparing Fick CO to LV outflow can reward the wrong physiology
\end{itemize}

\subsection{Treated the systemic quartet as one coupled physiological target}

The model now treats these as linked rather than fully independent:

\begin{verbatim}
SAP_mean, RAP_mean, CO_Lmin, SVR
\end{verbatim}

Reason:

\begin{itemize}[leftmargin=1.5em]
    \item \texttt{SVR = (SAP\_mean - RAP\_mean) / Qs}
    \item If these are optimized independently, the solver can improve pressure
    by suppressing flow or inflate flow while creating unrealistic resistance
\end{itemize}

\subsection{Reduced pulmonary waveform pressure matching for Reyna}

For Reyna, the active target now keeps:

\begin{itemize}[leftmargin=1.5em]
    \item \texttt{PAP\_mean}
    \item \texttt{QpQs}
\end{itemize}

and intentionally de-emphasizes:

\begin{itemize}[leftmargin=1.5em]
    \item \texttt{PAP\_min}
    \item \texttt{PAP\_max}
\end{itemize}

Reason:

\begin{itemize}[leftmargin=1.5em]
    \item The remaining bottleneck was not pulmonary polishing
    \item The model was already better at matching shunt and pulmonary mean
    pressure than at producing correct systemic flow and chamber-volume balance
\end{itemize}

\subsection{Added targeted systemic flow and volume guardrails for Reyna}

The Reyna-specific systemic-flow profile now:

\begin{itemize}[leftmargin=1.5em]
    \item strengthens \texttt{CO\_Lmin}, \texttt{SAP\_mean}, and \texttt{RAP\_mean}
    \item adds stronger penalty against high \texttt{RVEDV}
    \item adds stronger penalty against low \texttt{LVESV}
    \item adds stronger penalty against overly high \texttt{LVEF}
    \item keeps \texttt{CO\_Lmin} in the preferred primary set
\end{itemize}

This was implemented in:

\begin{itemize}[leftmargin=1.5em]
    \item \texttt{src/calibration/build\_case\_calibration\_profile.m}
    \item \texttt{src/calibration/objective\_calibration.m}
    \item \texttt{src/utils/select\_primary\_metrics.m}
\end{itemize}

\subsection{Added a regression test for the Reyna profile}

The Reyna-specific assumptions are now checked in:

\begin{itemize}[leftmargin=1.5em]
    \item \texttt{tests/test\_reyna\_systemic\_flow\_profile.m}
\end{itemize}

The test verifies:

\begin{itemize}[leftmargin=1.5em]
    \item raw protocol anthropometry remains active
    \item \texttt{PAP\_sys} and \texttt{PAP\_dia} stay omitted
    \item primary metrics include \texttt{CO\_Lmin}
    \item the volume-flow guard is active
\end{itemize}

\section{How \texttt{lb} and \texttt{ub} Are Defined}

The current calibration bounds are no longer arbitrary wide boxes. They are
defined centrally in:

\begin{itemize}[leftmargin=1.5em]
    \item \texttt{config/build\_parameter\_registry.m}
\end{itemize}

\subsection{General rule}

Each calibratable parameter gets:

\begin{itemize}[leftmargin=1.5em]
    \item a \texttt{baseline\_scaled} value after pediatric scaling
    \item a \texttt{seeded\_value} after clinical seeding
    \item a \texttt{bound\_anchor}
    \item an \texttt{lb}
    \item an \texttt{ub}
    \item a \texttt{bound\_source}
    \item a \texttt{confidence\_level}
\end{itemize}

This makes every active calibration knob traceable.

\subsection{Bound families}

\subsubsection{Coupled vascular resistance scales}

Examples:

\begin{itemize}[leftmargin=1.5em]
    \item \texttt{group.R\_sys\_scale}
    \item \texttt{group.R\_pul\_scale}
\end{itemize}

Default form:

\begin{verbatim}
lb = 0.25
ub = 2.80
\end{verbatim}

Then profile-level tightening can make the lower bound stricter. For Reyna:

\begin{itemize}[leftmargin=1.5em]
    \item \texttt{group.R\_sys\_scale}: \texttt{0.25} to \texttt{2.80}
    \item \texttt{group.R\_pul\_scale}: \texttt{0.45} to \texttt{2.80}
\end{itemize}

Source tag:

\begin{itemize}[leftmargin=1.5em]
    \item \texttt{kung2013\_bsa\_rc\_coupled\_resistance\_prior}
\end{itemize}

\subsubsection{VSD orifice coefficient}

Example:

\begin{itemize}[leftmargin=1.5em]
    \item \texttt{vsd.Cd}
\end{itemize}

Form:

\begin{verbatim}
lb = 0.20
ub = 1.20
\end{verbatim}

Source tag:

\begin{itemize}[leftmargin=1.5em]
    \item \texttt{orifice\_discharge\_coefficient\_literature\_plus\_expert\_box}
\end{itemize}

\subsubsection{Arterial compliances}

Examples:

\begin{itemize}[leftmargin=1.5em]
    \item \texttt{C.SAR}
    \item \texttt{C.PAR}
\end{itemize}

These are anchored to pulse-pressure estimates, not broad multiplier boxes.

For the Reyna active run:

\begin{itemize}[leftmargin=1.5em]
    \item \texttt{C.SAR}: \texttt{0.3785} to \texttt{0.6813}
\end{itemize}

Source tag:

\begin{itemize}[leftmargin=1.5em]
    \item \texttt{windkessel\_stroke\_volume\_over\_pulse\_pressure\_anchor}
\end{itemize}

\subsubsection{Resistances}

Example:

\begin{itemize}[leftmargin=1.5em]
    \item \texttt{R.SVEN}
\end{itemize}

Form:

\begin{verbatim}
lb = lower_mult * baseline_scaled
ub = upper_mult * baseline_scaled
\end{verbatim}

With venous resistances given slightly wider upper room.

For the Reyna active run:

\begin{itemize}[leftmargin=1.5em]
    \item \texttt{R.SVEN}: \texttt{0.0588} to \texttt{0.4119}
\end{itemize}

Source tag:

\begin{itemize}[leftmargin=1.5em]
    \item \texttt{kung2013\_sharp2000\_scaled\_resistance\_prior}
\end{itemize}

\subsubsection{Elastances}

Examples:

\begin{itemize}[leftmargin=1.5em]
    \item \texttt{E.LV.EA}
    \item \texttt{E.LV.EB}
    \item \texttt{E.RV.EA}
    \item \texttt{E.RV.EB}
\end{itemize}

These are anchored to pediatric elastance priors, then widened if needed to
include the clinically seeded start.

For the Reyna active run:

\begin{itemize}[leftmargin=1.5em]
    \item \texttt{E.LV.EA}: \texttt{7.0026} to \texttt{30.7613}
    \item \texttt{E.LV.EB}: \texttt{0.0840} to \texttt{0.3501}
    \item \texttt{E.RV.EA}: \texttt{1.3531} to \texttt{5.9046}
    \item \texttt{E.RV.EB}: \texttt{0.1041} to \texttt{0.4920}
\end{itemize}

Source tag:

\begin{itemize}[leftmargin=1.5em]
    \item \texttt{zhang2019\_pediatric\_elastance\_anchor}
\end{itemize}

\subsubsection{Unstressed volumes}

Examples:

\begin{itemize}[leftmargin=1.5em]
    \item \texttt{V0.LV}
    \item \texttt{V0.RV}
\end{itemize}

These are not free-floating broad search terms anymore. They are anchored to
blood-volume and preload consistency.

For the Reyna active run:

\begin{itemize}[leftmargin=1.5em]
    \item \texttt{V0.LV}: \texttt{3.8237} to \texttt{7.1376}
    \item \texttt{V0.RV}: \texttt{9.0034} to \texttt{18.3538}
\end{itemize}

Source tag:

\begin{itemize}[leftmargin=1.5em]
    \item \texttt{blood\_volume\_preload\_consistency\_prior}
\end{itemize}

\subsection{Why this matters}

The model now has two protections:

\begin{enumerate}[leftmargin=1.5em]
    \item The optimizer cannot drift into physiologically absurd values too easily.
    \item We can explain afterward whether a good fit was achieved cleanly or only
    by parking near the edge of the allowable space.
\end{enumerate}

\section{Reyna Active Bounds and Best Candidate}

Source artifacts:

\begin{itemize}[leftmargin=1.5em]
    \item \texttt{results/runs/20260511\_223816\_reyna\_pre\_surgery/tables/parameter\_registry\_pre\_surgery.csv}
    \item \texttt{results/runs/20260511\_223816\_reyna\_pre\_surgery/tables/parameter\_plausibility\_best\_candidate\_pre\_surgery.csv}
\end{itemize}

\begin{longtable}{lrrrrrl}
\toprule
Parameter & BaselineScaled & LB & UB & Best fit & Ratio & Plausibility \\
\midrule
\endhead
group.R\_sys\_scale & 1.0000 & 0.2500 & 2.8000 & 0.4016 & 0.4016 & WARNING \\
R.SVEN & 0.1471 & 0.0588 & 0.4119 & 0.3255 & 2.2126 & OK \\
group.R\_pul\_scale & 1.0000 & 0.4500 & 2.8000 & 0.5447 & 0.5447 & WARNING \\
C.SAR & 0.4079 & 0.3785 & 0.6813 & 0.5067 & 1.2423 & OK \\
E.LV.EA & 13.9824 & 7.0026 & 30.7613 & 7.0807 & 0.5064 & WARNING \\
E.LV.EB & 0.1400 & 0.0840 & 0.3501 & 0.2415 & 1.7246 & OK \\
E.RV.EA & 2.4602 & 1.3531 & 5.9046 & 5.6594 & 2.3004 & WARNING \\
E.RV.EB & 0.1892 & 0.1041 & 0.4920 & 0.1547 & 0.8174 & OK \\
V0.LV & 5.0983 & 3.8237 & 7.1376 & 6.0467 & 1.1860 & OK \\
V0.RV & 13.5954 & 9.0034 & 18.3538 & 9.1245 & 0.6711 & WARNING \\
\bottomrule
\end{longtable}

Interpretation:

\begin{itemize}[leftmargin=1.5em]
    \item The best Reyna scientific candidate stays inside all active bounds
    \item Several active parameters sit close to one edge of the allowed range
    \item That is why the run is scientifically informative but not yet clean
    enough to claim a finished patient-specific calibration
\end{itemize}

\section{Run Results So Far}

\subsection{Reyna sequence}

\begin{longtable}{p{4.0cm}p{4.7cm}rrp{3.1cm}}
\toprule
Run folder & Main change & Baseline RMSE & Best scientific RMSE & Status \\
\midrule
\endhead
\texttt{20260510\_225025\_reyna\_pre\_surgery} & first strong systemic-flow scientific candidate & 0.2513 & 0.1195 & REJECT, rollback \\
\texttt{20260511\_134027\_reyna\_pre\_surgery} & corrected-BSA experiment (\texttt{14.0 kg}, \texttt{98 cm}, \texttt{0.6173 m\^{}2}) & 0.2454 & 0.1087 & REJECT, rollback \\
\texttt{20260511\_203930\_reyna\_pre\_surgery} & raw PDF anthropometry + \texttt{PAP\_mean} only & 0.2699 & 0.1702 & REJECT, rollback \\
\texttt{20260511\_223816\_reyna\_pre\_surgery} & targeted systemic-flow polish & 0.2699 & 0.1287 & REJECT, rollback \\
\bottomrule
\end{longtable}

Latest Reyna best-candidate bottleneck metrics:

\begin{tabularx}{\textwidth}{lrrrr}
\toprule
Metric & Clinical & Best candidate & Error \% \\
\midrule
\texttt{RAP\_mean} & 5.000 & 5.343 & 6.86 \\
\texttt{PAP\_mean} & 15.000 & 16.049 & 6.99 \\
\texttt{SAP\_mean} & 71.300 & 63.559 & -10.86 \\
\texttt{QpQs} & 1.194 & 1.140 & -4.52 \\
\texttt{SVR} & 19.370 & 20.784 & 7.30 \\
\texttt{CO\_Lmin} & 3.423 & 2.801 & -18.17 \\
\texttt{RVEDV} & 30.500 & 38.079 & 24.85 \\
\texttt{LVESV} & 19.300 & 16.464 & -14.69 \\
\texttt{LVEF} & 0.528 & 0.610 & 15.45 \\
\bottomrule
\end{tabularx}

Interpretation:

\begin{itemize}[leftmargin=1.5em]
    \item The model is better than before at systemic consistency
    \item The remaining problem is still low effective systemic flow with RV
    volume overshoot and high LV ejection bias
    \item This is no longer mainly a pulmonary waveform problem
\end{itemize}

\subsection{Profile A}

Run folder:

\begin{itemize}[leftmargin=1.5em]
    \item \texttt{20260512\_070804\_patient\_profile\_a\_pre\_surgery}
\end{itemize}

Result summary:

\begin{tabular}{lr}
\toprule
Quantity & Value \\
\midrule
Baseline RMSE & 0.2444 \\
Best candidate RMSE & 0.2951 \\
Status & REJECT, rollback \\
\bottomrule
\end{tabular}

Main best-candidate errors:

\begin{tabular}{lr}
\toprule
Metric & Error \% \\
\midrule
\texttt{QpQs} & -62.52 \\
\texttt{SVR} & -52.36 \\
\texttt{CO\_Lmin} & +52.74 \\
\texttt{LVEDV} & -38.44 \\
\texttt{RVEDV} & -26.50 \\
\texttt{LVEF} & -22.29 \\
\bottomrule
\end{tabular}

Interpretation:

\begin{itemize}[leftmargin=1.5em]
    \item Calibration made A worse, not better
    \item This benchmark is not a good judge of the Reyna improvement strategy
    \item Profile A likely reflects benchmark inconsistency more than a clean
    model failure
\end{itemize}

\subsection{Profile B}

Run folder:

\begin{itemize}[leftmargin=1.5em]
    \item \texttt{20260512\_072450\_patient\_profile\_b\_pre\_surgery}
\end{itemize}

Result summary:

\begin{tabular}{lr}
\toprule
Quantity & Value \\
\midrule
Baseline RMSE & 0.2093 \\
Best candidate RMSE & 0.1880 \\
Status & REJECT, rollback \\
\bottomrule
\end{tabular}

Main best-candidate errors:

\begin{tabular}{lr}
\toprule
Metric & Error \% \\
\midrule
\texttt{RAP\_mean} & +38.67 \\
\texttt{CO\_Lmin} & +38.86 \\
\texttt{SVR} & -23.10 \\
\texttt{QpQs} & +18.92 \\
\texttt{PVR} & -13.70 \\
\texttt{SAP\_mean} & +10.11 \\
\bottomrule
\end{tabular}

Interpretation:

\begin{itemize}[leftmargin=1.5em]
    \item B improves slightly, but still rejects
    \item The residual pattern is different from Reyna
    \item This again argues against promoting the Reyna-specific polish blindly
    to all full-data cases
\end{itemize}

\subsection{Razka}

Run folder:

\begin{itemize}[leftmargin=1.5em]
    \item \texttt{20260512\_073515\_razka\_pre\_surgery}
\end{itemize}

Result summary:

\begin{tabular}{lr}
\toprule
Quantity & Value \\
\midrule
Baseline RMSE & 0.2290 \\
Best candidate RMSE & 0.1346 \\
Status & PROMISING\_NEAR\_MISS \\
Rollback applied & No \\
\bottomrule
\end{tabular}

Main best-candidate errors:

\begin{tabular}{lr}
\toprule
Metric & Error \% \\
\midrule
\texttt{RAP\_mean} & +3.64 \\
\texttt{PAP\_min} & -14.92 \\
\texttt{PAP\_mean} & -15.82 \\
\texttt{SAP\_mean} & +11.63 \\
\texttt{QpQs} & -1.57 \\
\bottomrule
\end{tabular}

Interpretation:

\begin{itemize}[leftmargin=1.5em]
    \item Razka is the strongest non-Reyna result so far
    \item It behaves like a real sparse-cath case, not like a synthetic benchmark
    \item Because Razka has no clinical CO or ventricular volume targets, it
    should be improved through the sparse-data path, not through Reyna-style
    systemic flow-volume tuning
\end{itemize}

\section{Current Scientific Position}

The work so far supports the following conclusion.

\subsection{What is working}

\begin{itemize}[leftmargin=1.5em]
    \item The calibration framework is much more defensible than before
    \item \texttt{lb} and \texttt{ub} are now traceable to clinical seeding and
    literature-informed priors
    \item Reyna improved meaningfully once CO was mapped to \texttt{Qs} and
    systemic variables were treated as one coupled target
    \item Razka shows that the real-data pipeline can produce a near-miss result
    without rollback
\end{itemize}

\subsection{What is not solved yet}

\begin{itemize}[leftmargin=1.5em]
    \item Reyna still undershoots effective systemic flow
    \item Reyna still overshoots \texttt{RVEDV}
    \item Some good Reyna fits still lean on boundary-adjacent parameters
    \item Synthetic A and B are not reliable enough to validate the same fix path
\end{itemize}

\subsection{Best next interpretation}

\begin{itemize}[leftmargin=1.5em]
    \item Reyna remains the best case for systemic preload-flow-volume cleanup
    \item Razka should be used to improve the sparse real-data path
    \item Profile A and B should remain synthetic stress benchmarks, not core
    scientific validation cases
\end{itemize}

\section{Files Most Relevant to This Progress}

\begin{itemize}[leftmargin=1.5em]
    \item \texttt{config/patient\_reyna.m}
    \item \texttt{docs/clinical\_data\_dictionary.md}
    \item \texttt{docs/reyna\_pak\_dipo\_systemic\_flow\_revision.md}
    \item \texttt{config/build\_parameter\_registry.m}
    \item \texttt{src/calibration/build\_case\_calibration\_profile.m}
    \item \texttt{src/calibration/objective\_calibration.m}
    \item \texttt{tests/test\_reyna\_systemic\_flow\_profile.m}
\end{itemize}

\end{document}
